\documentclass[11pt, letterpaper]{article}

% ----- Packages -----
\usepackage[utf8]{inputenc}
\usepackage{amsmath, amssymb, amsthm, mathrsfs, amsfonts, mathtools}
\usepackage{geometry}
\geometry{margin=1.15in}
\usepackage{hyperref}
\usepackage{xcolor}
\usepackage{titlesec}
\usepackage{enumitem}
\usepackage{authblk}
\usepackage{booktabs}
\usepackage{graphicx}

% ----- Formatting -----
\hypersetup{
    colorlinks=true,
    linkcolor=blue,
    filecolor=magenta,
    urlcolor=blue,
    citecolor=black,
}
\titleformat{\section}{\large\bfseries}{\thesection.}{1em}{}
\titleformat{\subsection}{\normalsize\bfseries}{\thesubsection}{1em}{}

% ----- Math Commands -----
\newcommand{\R}{\mathbb{R}}
\newcommand{\C}{\mathbb{C}}
\newcommand{\Z}{\mathbb{Z}}
\newcommand{\G}{G^*}
\DeclareMathOperator{\AGM}{AGM}

% ----- Theorem Environments -----
\newtheorem{theorem}{Theorem}[section]
\newtheorem{lemma}[theorem]{Lemma}
\theoremstyle{definition}
\newtheorem{definition}{Definition}[section]
\theoremstyle{remark}
\newtheorem{remark}{Remark}[section]

% ----- Epistemic Tags -----
\newcommand{\epi}[1]{\textbf{\textsc{[#1]}}}

% ----- Document Info -----
\title{\vspace{-3em}\textbf{The Ontic Mathematical Foundations: From Counting to Coupling\\ \large The Constant Chain $\gamma \to \varpi \to M \to \pi \to \G$}}
\author[1]{William J. Steinmetz III}
\affil[1]{\textit{Foundational Ternary Dynamics Research}}
\date{February 14, 2026 \\ \vspace{0.5em} \small Framework: FTD v5.24}

\begin{document}

\maketitle

\begin{abstract}
We formalize the ontic foundation of mathematical constants in Foundational Ternary Dynamics (FTD) as a directed chain of emergence: $\gamma \to \varpi \to M \to \pi \to \G$, where $\gamma$ is the Euler--Mascheroni constant, $\varpi$ is the lemniscate constant, $M = \AGM(1, \sqrt{2})$ is the Gauss arithmetic-geometric mean, $\pi$ is Archimedes' constant, and $\G$ is the lemniscate-alpha constant of FTD. The chain is grounded by a rigorous structural claim: \textbf{$\gamma$ is the inversion term where the discrete integer world first makes contact with the continuous analytical world, and this contact is logarithmic in character.} Without $\gamma$, there is no bridge from discrete counting to the elliptic and lemniscatic structures from which the fine-structure constant ($\alpha$) emerges. The ordering is defined not by numerical value, but by ontological depth.
\end{abstract}

\vspace{1em}
\hrule
\vspace{1.5em}

\noindent \textit{\textbf{Epistemic Tags:} Throughout this document, we utilize strict epistemic bounding: \epi{Theorem} for structural results, \epi{Definition} for defining relations, \epi{Standard} for classical analysis, \epi{Selection} for ordering, and \epi{Conjecture} for physical interpretation.}

\vspace{1em}

%===============================================================================
\section{The Five Constants}
%===============================================================================

\subsection{The Euler--Mascheroni Constant $\gamma$}

\begin{equation}
    \gamma = \lim_{n \to \infty} \left( \sum_{k=1}^{n} \frac{1}{k} - \ln n \right) = 0.5772156649\dots
\end{equation}

\noindent\textbf{What it requires:} Only the positive integers and the concept of a limit.

\noindent\textbf{What it encodes:} The exact offset between \textit{discrete counting} (the harmonic series $H_n = 1 + 1/2 + 1/3 + \cdots + 1/n$) and \textit{continuous scaling} (the natural logarithm $\ln n$). This makes $\gamma$ the \textbf{universal conversion constant} between the arithmetic world of integers and the analytic world of continuous functions.

\noindent\textbf{The inversion property:} Rearranging $H_n \approx \ln(n) + \gamma$, we find
\begin{equation}
    \exp(H_n - \gamma) \approx n
\end{equation}
Without $\gamma$, one cannot invert from continuous logarithmic scaling back to discrete counting. It is the constant that renders the integers recoverable from their logarithmic shadow.

\noindent\textbf{Why $\gamma$ is ontologically first:} Every subsequent constant in the chain depends on the Gamma function $\Gamma(z)$, and $\gamma$ appears in its Weierstrass product representation:
\begin{equation}
    \frac{1}{\Gamma(z)} = z \cdot e^{\gamma z} \cdot \prod_{n=1}^{\infty}\left[\left(1 + \frac{z}{n}\right) e^{-z/n}\right]
\end{equation}
Here $\gamma$ is the \textbf{exponential rate} that controls how the Gamma function scales. It is embedded multiplicatively---as $\exp(\gamma z)$---into every evaluation of $\Gamma$, including the critical value $\Gamma(1/4)$ from which $\varpi$ and $\G$ are built.

\subsection{The Lemniscate Constant $\varpi$}

\begin{equation}
    \varpi = \frac{\Gamma(1/4)^2}{2\sqrt{2\pi}} = 2.6220575543\dots
\end{equation}

\noindent\textbf{What it requires:} $\gamma$ (through $\Gamma(1/4)$), the integer 4, and $\pi$.

\noindent\textbf{What it encodes:} The half-period of the Bernoulli lemniscate $r^2 = \cos(2\theta)$---the ``$\pi$ of the lemniscate.'' Just as $\pi$ measures the circle, $\varpi$ measures the figure-eight: the first self-crossing algebraic curve, encoding self-reference geometrically.

\noindent\textbf{Connection to $\gamma$:} The digamma function satisfies
\begin{equation}
    \psi(1/4) = -\gamma - \frac{\pi}{2} - 3\ln 2
\end{equation}
This identity shows that $\gamma$ is \textit{analytically present} in the logarithmic derivative of $\Gamma$ at the critical point $z = 1/4$. Since $\Gamma(1/4)$ determines $\varpi$, $\gamma$ flows through to $\varpi$ via the derivative chain:
\begin{equation}
    \gamma \;\xrightarrow{\text{Weierstrass}}\; \Gamma(1/4) \;\xrightarrow{\text{definition}}\; \varpi
\end{equation}

\subsection{The Gauss Constant $M$}

\begin{equation}
    M = \AGM(1, \sqrt{2}) = 1.1981402347\dots
\end{equation}

\noindent\textbf{What it requires:} The arithmetic-geometric mean iteration applied to 1 and $\sqrt{2}$.

\noindent\textbf{What it encodes:} The rate of convergence of the AGM process---itself a bridge between arithmetic means (discrete averaging) and geometric means (multiplicative scaling). The AGM is the fastest-converging classical algorithm, doubling precision with each step.

\epi{Theorem} \textbf{Exact relationship:}
\begin{equation}
    \varpi = \frac{\pi}{M} \qquad \Longleftrightarrow \qquad M = \frac{\pi}{\varpi}
\end{equation}

This identity reveals $M$ as the \textbf{conversion factor between circular geometry ($\pi$) and lemniscatic geometry ($\varpi$)}. The reciprocal $1/M \approx 0.8346$ is the classical Gauss constant.

\noindent\textbf{Note on ordering:} $M$ sits between $\varpi$ and $\pi$ in the chain because it is the bridge that \textbf{converts} between them. $\varpi$ is more fundamental (requiring only $\Gamma(1/4)$ and hence only $\gamma$ and 4); $M$ requires the AGM concept; $\pi$ requires circular geometry or, equivalently, the combination $M \cdot \varpi$.

\subsection{Archimedes' Constant $\pi$}

\begin{equation}
    \pi = M \cdot \varpi = \frac{4\varpi^2}{\G^{2}} = 3.1415926536\dots
\end{equation}

\noindent\textbf{What it requires:} The circle, or equivalently, the product of $\varpi$ and $M$.

\noindent\textbf{What it encodes:} The ratio of circumference to diameter---the constant of \textbf{circular closure}. In FTD, $\pi$ is not ontologically primitive. It is the product of two more fundamental constants:
\begin{equation}
    \pi = \varpi \cdot M = \varpi \cdot \AGM(1, \sqrt{2})
\end{equation}

This decomposition reveals $\pi$ as a \textbf{composite}: the lemniscatic half-period ($\varpi$) scaled by the arithmetic-geometric bridge ($M$).

\epi{Selection} The claim that $\pi$ is ontologically derivative of $\varpi$ is a selection principle. Mathematically, any of $\{\gamma, \varpi, M, \pi, \G\}$ can be expressed in terms of the others. The ordering reflects the principle that constants requiring less structure are more fundamental.

\subsection{The Lemniscate-Alpha Constant $\G$}

\begin{equation}
    \G = \frac{\sqrt{2}\;\Gamma(1/4)^2}{2\pi} = \frac{2\varpi}{\sqrt{\pi}} = 2.9586751192\dots
\end{equation}

\noindent\textbf{What it requires:} All four preceding constants, or equivalently, $\varpi$ and $\pi$ together.

\noindent\textbf{What it encodes:} The master coefficient of the FTD quadratic \cite{Steinmetz2026_Zenodo}:
\begin{equation}
    x^2 - 16\,\G^{2}\,x + 16\,\G^{3} = 0
\end{equation}
whose roots are $x_+ = 137.036\ldots \approx 1/\alpha$ and $x_- = 3.024\ldots \approx N_c$.

\noindent\textbf{Why $\G$ is ontologically last:} It is the most ``constructed'' of the five. It requires $\sqrt{2}$ (from the Gauss constraint geometry), $\Gamma(1/4)^2$ (which carries $\gamma$ through Weierstrass), and $\pi$ (which carries $M$ and $\varpi$). $\G$ is where all the preceding structure \textbf{converges} into a single number that, through the master quadratic, produces physics.

%===============================================================================
\section{The Chain as Ontological Ordering}
%===============================================================================

\subsection{The Principle}

The ordering $\gamma \to \varpi \to M \to \pi \to \G$ is not by numerical magnitude but by \textbf{ontological depth}: the amount of mathematical structure required for each constant's definition.

\begin{table}[h]
\centering
\begin{tabular}{@{}clcll@{}}
\toprule
\textbf{Level} & \textbf{Constant} & \textbf{Value} & \textbf{Requires} & \textbf{Encodes} \\
\midrule
0 & $\gamma$ & 0.5772\ldots & Integers + limit & Discrete $\leftrightarrow$ continuous bridge \\
1 & $\varpi$ & 2.6221\ldots & $\gamma$ (via $\Gamma$), integer 4 & Self-crossing geometry \\
2 & $M$ & 1.1981\ldots & AGM iteration & Arithmetic $\leftrightarrow$ geometric bridge \\
3 & $\pi$ & 3.1416\ldots & $\varpi \times M$ & Circular closure \\
4 & $\G$ & 2.9587\ldots & $\sqrt{2}$, $\varpi$, $\pi$ & Master quadratic $\to \alpha$ \\
\bottomrule
\end{tabular}
\caption{The ontic hierarchy of mathematical constants.}
\end{table}

\subsection{Each Level Adds Structure}

\textbf{Level 0 $\to$ 1 ($\gamma \to \varpi$):} From discrete/continuous bridging to self-crossing geometry. The Weierstrass product converts $\gamma$ (a limit of sums) into $\Gamma(1/4)$ (a special function value), and the definition $\varpi = \Gamma(1/4)^2/(2\sqrt{2\pi})$ crystallizes this into the period of a self-referential curve. This step adds \textbf{geometric self-reference}.

\textbf{Level 1 $\to$ 2 ($\varpi \to M$):} From lemniscatic period to the AGM convergence rate. $M = \pi/\varpi$ is defined by the arithmetic-geometric mean, which iterates between two types of averaging. This step adds \textbf{iterative convergence}.

\textbf{Level 2 $\to$ 3 ($M \to \pi$):} From AGM rate to circular constant. $\pi = \varpi \cdot M$ combines the self-crossing period with the convergence bridge. This step adds \textbf{closure} (the circle is the simplest closed curve).

\textbf{Level 3 $\to$ 4 ($\pi \to \G$):} From circular to lemniscatic with $\sqrt{2}$ scaling. $\G = 2\varpi/\sqrt{\pi} = \sqrt{2} \cdot \Gamma(1/4)^2/(2\pi)$ introduces the $\sqrt{2}$ from the FCC/Gauss constraint geometry. This step adds \textbf{constraint physics}---the discrete lattice's Gauss law.

\subsection{$\gamma$ as the Logarithmic Inversion Term}

The deepest claim: \textbf{$\gamma$ is where numbers first scale logarithmically into the decimal space.}

The harmonic series $H_n = 1 + 1/2 + \cdots + 1/n$ grows as $\ln(n) + \gamma$. This means the ``natural'' way to accumulate the reciprocals of integers (a purely arithmetic operation on $1, 2, 3, \ldots$) produces a logarithmic function plus a correction term $\gamma$. The integers, through their reciprocal sum, \textit{generate} logarithmic scaling---and $\gamma$ is the remainder.

More precisely:
\begin{itemize}[noitemsep]
    \item \textbf{Discrete world:} $H_n = \sum 1/k$ (sum over integers)
    \item \textbf{Continuous world:} $\ln(n)$ (integral $\int dx/x$)
    \item \textbf{Bridge:} $\gamma = H_n - \ln(n)$ in the limit
\end{itemize}

The exponential form $e^{-\gamma} \approx 0.5615$ acts as a \textbf{decimal discount factor}: it governs how the discrete harmonic world undershoots the continuous logarithmic world. By Mertens' theorem \cite{Mertens1874}, $e^{-\gamma}$ even controls the distribution of primes:
\begin{equation}
    \prod_{p \leq N,\; p\text{ prime}} \left(1 - \frac{1}{p}\right) \;\sim\; \frac{e^{-\gamma}}{\ln N}
\end{equation}
So $\gamma$ is not merely a correction term---it is the constant that governs how \textbf{prime factorization} (the most fundamental discrete structure) maps to \textbf{logarithmic density} (the most fundamental continuous measure).

%===============================================================================
\section{The Complete Derivation Chain}
%===============================================================================

\subsection{From $\gamma$ to $\alpha$}

The full chain, with each step justified:

\begin{equation}
    \boxed{\gamma \;\xrightarrow{\text{Weierstrass}}\; \Gamma(1/4) \;\xrightarrow{\text{def.}}\; \varpi \;\xrightarrow{\text{AGM}}\; M \;\xrightarrow{\times\,\varpi}\; \pi \;\xrightarrow{\sqrt{2}\,\text{scaling}}\; \G \;\xrightarrow{\text{quadratic}}\; \alpha}
\end{equation}

\begin{enumerate}[noitemsep]
    \item $\gamma = 0.5772\ldots$ (integers + limit)
    \item Weierstrass product: $1/\Gamma(z) = z \cdot e^{\gamma z} \cdot \prod\ldots$; evaluate at $z = 1/4$
    \item $\Gamma(1/4) = 3.6256\ldots$
    \item Definition: $\varpi = \Gamma(1/4)^2/(2\sqrt{2\pi}) = 2.6221\ldots$
    \item AGM identity: $M = \pi/\varpi = \AGM(1,\sqrt{2}) = 1.1981\ldots$
    \item Product: $\pi = \varpi \cdot M = 3.1416\ldots$
    \item Scaling: $\G = 2\varpi/\sqrt{\pi} = \sqrt{2}\cdot\Gamma(1/4)^2/(2\pi) = 2.9587\ldots$
    \item Master quadratic: $x^2 - 16\G^2 x + 16\G^3 = 0$
    \item $x_+ = 137.036\ldots \approx 1/\alpha$, \quad $x_- = 3.024\ldots \approx N_c$
\end{enumerate}

\subsection{The Defining Relations}

\epi{Definition} The entire chain rests on exactly \textbf{two independent definitions} plus one \textbf{FTD-specific claim}:

\begin{definition}[Lemniscate Constant]
$\displaystyle \varpi = \frac{\Gamma(1/4)^2}{2\sqrt{2\pi}}$
\end{definition}

\begin{definition}[Lemniscate-Alpha Constant]
$\displaystyle \G = \frac{\sqrt{2}\;\Gamma(1/4)^2}{2\pi} = \frac{2\varpi}{\sqrt{\pi}}$
\end{definition}

\epi{Conjecture} The master quadratic with coefficient 16 (from lattice degrees of freedom):
\begin{equation}
    x^2 - 16\,\G^{2}\,x + 16\,\G^{3} = 0
\end{equation}
has roots $x_+ = 137.036\ldots \approx 1/\alpha$ and $x_- = 3.024\ldots \approx N_c$.

\subsection{The $\gamma$--$\varpi$ Connection: How Counting Becomes Geometry}

\epi{Theorem} This is the genuinely meaningful part of the chain---the mechanism by which the discrete-to-continuous bridge constant $\gamma$ flows into the lemniscatic structure.

The \textbf{Weierstrass product} provides the link:
\begin{equation}
    \frac{1}{\Gamma(z)} = z \cdot e^{\gamma z} \cdot \prod_{n=1}^{\infty}\left[\left(1 + \frac{z}{n}\right) e^{-z/n}\right]
\end{equation}
Evaluated at $z = 1/4$ (the FTD-selected evaluation point), $\gamma$ enters as the factor $\exp(\gamma/4)$ inside $\Gamma(1/4)$. Since $\varpi$ and $\G$ depend on $\Gamma(1/4)^2$, $\gamma$ propagates through as $\exp(\gamma/2)$. This is not a numerical coincidence---it is a \textbf{structural consequence} of how the Gamma function is built from the integers.

The digamma function (logarithmic derivative of $\Gamma$) makes this connection explicit:
\begin{align}
    \psi(1/4) &= -\gamma - \frac{\pi}{2} - 3\ln 2 &\text{\epi{Standard}---Gauss digamma theorem} \\
    \psi(1/3) &= -\gamma - \frac{3}{2}\ln 3 - \frac{\pi}{2\sqrt{3}} &\text{\epi{Standard}---Gauss digamma theorem}
\end{align}
These are classical identities due to Gauss (1813) \cite{Gauss1813}. What is structurally meaningful for FTD is that the evaluation points $z = 1/3$ and $z = 1/4$ correspond to the two primary framework integers $N_c = 3$ and $N_{\text{base}} = 4$. Setting both expressions equal (since both define $\gamma$) yields:
\begin{equation}
    \psi(1/4) - \psi(1/3) = \frac{3}{2}\ln 3 + \frac{\pi}{2\sqrt{3}} - \frac{\pi}{2} - 3\ln 2
\end{equation}
This is a non-trivial \textbf{cross-constraint} linking the digamma at the two FTD integers through a specific combination of $\pi$, $\ln 2$, $\ln 3$, and $\sqrt{3}$.

\subsection{Standard Results Used}

\epi{Standard} The chain uses several classical results from analysis. These are \textbf{tools}, not FTD discoveries:

\begin{table}[h]
\centering
\small
\begin{tabular}{@{}lll@{}}
\toprule
\textbf{Result} & \textbf{Source} & \textbf{Role in Chain} \\
\midrule
$\Gamma(z)\cdot\Gamma(1-z) = \pi/\sin(\pi z)$ & Euler reflection & Relates $\Gamma(1/4)$ to $\Gamma(3/4)$ \\
Gauss multiplication formula & Gauss (1812) & Connects $\Gamma$ at multiples of $1/n$ \\
$\psi(1/n)$ closed forms & Gauss (1813) & Links $\gamma$ to $\Gamma(1/4)$ analytically \\
$\varpi \cdot M = \pi$ & AGM--elliptic identity & Defines $M$ as bridge \\
$\zeta(2n) = (-1)^{n+1}B_{2n}(2\pi)^{2n}/(2\cdot(2n)!)$ & Euler (1735) & Even zeta from Bernoulli numbers \\
$\prod(1-1/p) \sim e^{-\gamma}/\ln N$ & Mertens (1874) & $\gamma$ governs prime distribution \\
\bottomrule
\end{tabular}
\caption{Classical results employed in the derivation chain.}
\end{table}

%===============================================================================
\section{Structural Theorems}
%===============================================================================

\subsection{The Minimal Generating Set}

\begin{theorem}[Minimal Generating Set]
The entire ontic chain is generated by exactly \textbf{two constants}: $\gamma$ (Euler--Mascheroni) and $\pi$ (Archimedes).
\end{theorem}

\textit{Proof sketch.} Given $\gamma$ and $\pi$:
\begin{enumerate}[noitemsep]
    \item $\Gamma(1/4)$ is determined by the Weierstrass product (requiring only $\gamma$ and the integer 4)
    \item $\varpi = \Gamma(1/4)^2/(2\sqrt{2\pi})$
    \item $M = \pi/\varpi = \AGM(1,\sqrt{2})$
    \item $\G = 2\varpi/\sqrt{\pi} = \sqrt{2}\cdot\Gamma(1/4)^2/(2\pi)$
    \item $x_+, x_-$ from the master quadratic
    \item $\alpha = 1/x_+$ (with precision corrections using framework integers)
\end{enumerate}

The circularity ($\pi$ appears in step 2, which seems to require $\pi$ before generating it) resolves because \textbf{both $\gamma$ and $\pi$ are independently definable from the integers alone}: $\gamma = \lim[H_n - \ln(n)]$ and $\pi = 4\arctan(1) = 4\sum(-1)^n/(2n+1)$.

\textbf{Implication:} The fine structure constant $\alpha$ is determined by exactly two transcendental constants from number theory, plus the framework integers $\{3, 4, 7, 13\}$.

\subsection{The $\exp(\gamma/2)$ Universal Scaling}

\begin{theorem}[Universal Rescaling]
The factor $\exp(\gamma/2)$ is the \textbf{universal rescaling} that removes $\gamma$ from all lemniscatic constants simultaneously.
\end{theorem}

Define the ``$\gamma$-free'' constants:
\begin{itemize}[noitemsep]
    \item $\Gamma_0(1/4) = \Gamma(1/4) \cdot \exp(\gamma/4)$
    \item $\varpi_0 = \varpi \cdot \exp(\gamma/2) = \Gamma_0(1/4)^2/(2\sqrt{2\pi})$
    \item $\G_0 = \G \cdot \exp(\gamma/2)$
\end{itemize}

\textbf{Why $\exp(\gamma/2)$?} Because $\gamma$ enters the chain \textbf{only} through $\Gamma(1/4)$ via the Weierstrass product, where it contributes a factor of $\exp(\gamma/4)$. Since $\varpi$ and $\G$ depend on $\Gamma(1/4)^2$, the $\gamma$-dependence always enters as $\exp(\gamma/4)^2 = \exp(\gamma/2)$.

\textbf{Physical reading:} The factor $\exp(-\gamma/2) \approx 0.749$ is the ``discount'' that the discrete-to-continuous bridge ($\gamma$) applies to the idealized lemniscatic constants. Without this discount, $\varpi_0 \approx 3.499$ and $\G_0 \approx 3.949$---both larger than their actual values. The integers, through their logarithmic shadow, \textbf{compress} the lemniscatic constants.

%===============================================================================
\section{Relation to the FTD Hierarchy}
%===============================================================================

\subsection{Where This Fits}

The existing FTD ontological hierarchy begins at Level 0 with the Pure Integral $I_4$. The ontic constant chain provides the \textbf{sub-structure beneath $I_4$}:

\begin{table}[h]
\centering
\begin{tabular}{@{}lll@{}}
\toprule
\textbf{Prior Hierarchy} & \textbf{This Document} & \textbf{Constant} \\
\midrule
Level $-3$: Absolute Void & --- & (no mathematics) \\
Level $-2$: Pregnant Void & --- & (potentiality) \\
Level $-1$: First Distinction & Level 0 & $\gamma$ (integers emerge) \\
Level 0: Self-Reference ($n=4$) & Level 0$\to$1 & $\gamma \to \Gamma(1/4)$ \\
Level 1: Pure Integral $I_4$ & Level 1 & $I_4 = \varpi/2$ \\
Level 5: Lemniscatic & Levels 1--4 & $\varpi \to M \to \pi \to \G$ \\
\bottomrule
\end{tabular}
\caption{Integration of the ontic chain with the FTD hierarchy.}
\end{table}

\noindent\textbf{Key identification:} $I_4 = \varpi/2$. The ``Pure Integral'' of the prior hierarchy is exactly half the lemniscate constant. The ontic chain $\gamma \to \varpi$ provides the foundation \textit{beneath} $I_4$ by showing where the Gamma function (and hence $\varpi$) gets its scaling from.

\subsection{The Role of the Integer 4}

In the prior hierarchy, $n = 4$ is selected by self-reference. In the ontic chain, 4 appears at the \textbf{evaluation point} $z = 1/4$ where the Weierstrass product (carrying $\gamma$) is evaluated to produce $\Gamma(1/4)$. The two hierarchies converge:

\begin{itemize}[noitemsep]
    \item \textbf{Prior:} Self-reference selects $n = 4$ $\to$ $I_4 = \int_0^1 dx/\sqrt{1-x^4}$
    \item \textbf{Ontic:} $\gamma$ flows through $\Gamma(z)$ at $z = 1/4$ $\to$ $\Gamma(1/4) \to \varpi = 2I_4$
\end{itemize}

Both paths arrive at the same destination: the lemniscate constant $\varpi$. The integer 4 is the \textbf{meeting point} of the self-referential argument (why $n = 4$) and the analytic argument (why $z = 1/4$ in the Weierstrass product).

%===============================================================================
\section{Summary and Physical Interpretation}
%===============================================================================

\subsection{The Ontic Chain}

\begin{equation}
    \boxed{\gamma \;\xrightarrow{\text{Weierstrass}}\; \Gamma(1/4) \;\xrightarrow{\text{def.}}\; \varpi \;\xrightarrow{\text{AGM}}\; M \;\xrightarrow{\times\,\varpi}\; \pi \;\xrightarrow{\sqrt{2}\,\text{scaling}}\; \G \;\xrightarrow{\text{quadratic}}\; \alpha}
\end{equation}

\subsection{Physical Interpretation}

\epi{Conjecture} $\gamma$ is the \textbf{first constant that exists} once the integers exist and the concept of a limit is available. It encodes the irreducible gap between counting and measuring---between the discrete and the continuous. Every subsequent constant in the chain inherits this gap through the Gamma function's Weierstrass product, which carries $\gamma$ as an exponential rate factor.

The physical implication: \textbf{the fine structure constant $\alpha$ inherits its specific numerical value from $\gamma$ through a chain of exact algebraic relationships.} The integers provide the skeleton ($\{3, 4, 7, 13\}$); $\gamma$ provides the flesh (the continuous scaling); and the lemniscatic structure provides the self-referential closure that binds them together.

\subsection{Epistemic Status}

\begin{table}[h]
\centering
\small
\begin{tabular}{@{}lll@{}}
\toprule
\textbf{Claim} & \textbf{Tag} & \textbf{Notes} \\
\midrule
$\gamma$ as logarithmic inversion term & \epi{Theorem} & Follows from $H_n = \ln(n) + \gamma + O(1/n)$ \\
Weierstrass: $\gamma$ flows into $\Gamma(1/4)$ & \epi{Theorem} & Standard analysis; structural mechanism \\
$\exp(\gamma/2)$ universal scaling & \epi{Theorem} & Consequence of $\Gamma(1/4)^2$ dependence \\
Minimal generating set $\{\gamma, \pi\}$ & \epi{Theorem} & Constructive proof \\
Digamma cross-constraint & \epi{Standard} & Gauss (1813); FTD integers have closed forms \\
Ordering by ontological depth & \epi{Selection} & Principled but not uniquely determined \\
Master quadratic yields $x_+ \approx 1/\alpha$ & \epi{Conjecture} & 1.26 ppm match; no mechanism derived \\
Coefficient 16 from lattice DoF & \epi{Selection} & Argued from geometry; product rule imposed \\
\bottomrule
\end{tabular}
\caption{Epistemic classification of all claims.}
\end{table}

\vspace{2em}
\hrule
\vspace{1.5em}

\begin{thebibliography}{99}

\bibitem{Steinmetz2026_Zenodo}
W. J. Steinmetz III,
``Gauge Couplings from Discrete Spacetime Geometry,''
\textit{Foundational Ternary Dynamics Research}, Zenodo (January 2026).
\url{https://zenodo.org/records/18227181}

\bibitem{Gauss1813}
C. F. Gauss,
``Disquisitiones generales circa seriem infinitam,''
\textit{Commentationes societatis regiae scientiarum Gottingensis recentiores}, \textbf{2}, 1--46 (1813).
\newblock \textit{[Closed-form evaluations of the digamma function at rational arguments.]}

\bibitem{Mertens1874}
F. Mertens,
``Ein Beitrag zur analytischen Zahlentheorie,''
\textit{Journal f\"ur die reine und angewandte Mathematik}, \textbf{78}, 46--62 (1874).
\newblock \textit{[Asymptotic product over primes: $\prod_{p \leq N}(1-1/p) \sim e^{-\gamma}/\ln N$.]}

\bibitem{Euler1735}
L. Euler,
``De summis serierum reciprocarum,''
\textit{Commentarii academiae scientiarum Petropolitanae}, \textbf{7}, 123--134 (1735).
\newblock \textit{[Solution to the Basel problem and evaluation of $\zeta(2n)$ via Bernoulli numbers.]}

\end{thebibliography}

\end{document}